%!TEX ROOT=sample-thesis.tex
%!TEX PROGRAM=xelatex

\chapter{Results and Discussions}

In this chapter, the results of the user case study are presented and discussed to evaluate the proposed design. The recorded user activities in think-aloud observation were examined and discussed to identify user cognition sequence and behavior. The findings from the user case study then validated based on experts’ interview results. The perceptual usefulness and effectiveness of the proposed design including each view panels were further validated. Furthermore, discussions on the findings based on design fulfilment, domain user behavior, and design limitations are presented.
\\ \\ \indent
The remainder of this chapter is as follows. Section 4.1 present and discuss the results of think-aloud observation on each view. Section 4.2 discuss on expert validation towards findings from user case study. In Section 4.3, significant findings from the user case study are further discussed. Finally, the conclusion of this chapter is presented in Section 4.4.

\section{User Case Study Results}
This section presents and discusses the user case study results from the think-aloud observation with domain expert. Domain experts’ experience in performing MOOC analysis using the proposed visual analytics design are reported. This study then further examines their analysis activity on each view and thought process during the observation.

\subsection{User Cognition on Overview View}
Experts’ cognition experience when using the Overview features of the proposed visual analytics are examined and described in Table 4.1. The results of cognition experience on each task are discussed as follows.

\begin{table}[hbt!]
\caption{Cognition experience of domain experts in completing tasks using Overview}
\centering
\begin{tabular}{ ccccccc }
\hline
\textbf{Task} 	& \textbf{DE1} 	& \textbf{DE2} 	& \textbf{DE3} 	& \textbf{DE4} 	& \textbf{DE5} 	& \textbf{DE6} 	\\ \hline
\multicolumn{7}{l}{\textbf{Identify}}	\\ \hline
$OI_1$			& Moderate 		& Easy 			& Easy 			& Easy 			& Easy 			& Easy 			\\ 
$OI_2$			& Easy 			& Easy 			& Easy 			& Easy 			& Moderate 		& Easy 			\\ 
$OI_3$			& Moderate 		& Easy 			& Easy 			& Easy 			& Easy 			& Moderate 		\\ 
%\multicolumn{7}{|l|}{} \\
%\multicolumn{7}{|l|}{\textbf{Find Extremum}} \\ \hline
$OE_1$	& Easy & Easy & Easy & Moderate & Easy & Moderate 	\\ 
$OE_2$	& Easy & Easy & Easy & Easy & Easy & Easy 	\\ 
$OE_3$	& Moderate & Easy & Easy & Easy & Easy & Easy \\
%\multicolumn{7}{|l|}{} \\
%\multicolumn{7}{|l|}{\textbf{Compare}} \hline
$OC_1$			& Moderate 		& Easy 			& Moderate 		& Moderate 		& Easy 			& Moderate 		\\ 
$OC_2$			& Moderate 		& Easy 			& Easy 			& Moderate 		& Easy 			& Moderate 		\\ 
$OC_3$			& Moderate 		& Easy 			& Easy 			& Easy 			& Easy 			& Easy 			\\
\hline
\end{tabular}
\end{table}

\subsubsection{Identify}
\textbf{$OI_1$: Identify total number of learners in the course} \\
Generally, all participant was able to accomplish this task. However, few participants require additional time to recognize the presented information. In this task scenario, it was found that the use of color, size, and positioning is crucial in rendering first-hand information. Expert 1 remarked \textit{“I browsed until the bottom. I thought that perhaps it looks like accounting statement, with great total”}. However, this study considered it as uncommon scenario since the other participants were able to identify the item from top. Moreover, this study learns that appropriate use of colors without compromising cognition attention is crucial in presenting overview information. Expert 4 expressed excitement towards the colors in the panel, but she mistaken the values in colored legend as total learners. Furthermore, Expert 1 commented \textit{“Since the total number of learners is important, if it was red-colored, I will definitely be able to recognize it. Not gray though.”}. These findings on colors utilization correlates with previous work that emphasizes on effective color mapping.
\\[\baselineskip]\indent
Time for some maths.

\begin{equation}
\left[M\frac{\partial }{\partial M}+\beta(g)\frac{\partial }{\partial g}+n\gamma\right] G^{(n)}(x_1,x_2,\ldots,x_n;M,g)=0,
\end{equation}
